%【译文使用说明】
% 完整中文译文逐句对应当前英文，仅用于源码审阅，不进入PDF。原有简短中文注释保留为写作说明或核验提醒；二者含义不同。
% ============================================================================
%【源码审阅导航｜第 6 章：结论与未来工作】
% 实际正文编号为第 6 章；文件名：sections/07_conclusion.tex。标记仅供源码审阅，不显示在PDF中。
% 正文按节编号；公式前后可分为编辑块。修改正文时保留引用标签、数学符号和图表数据；待补项沿用原有说明。
% ============================================================================
%【待补正文】本文件当前只有章节标题与写作提示，尚未写入结论正文。
% 第六章使用本文件；06_discussion.tex未被main.tex引用。
% ---------------------------------------------------------------------------
%【第 6 节｜Conclusion and Future Work】
% ---------------------------------------------------------------------------
\section{Conclusion and Future Work}
%【章节标题译文】
% 结论与未来工作
\label{sec:conclusion}

% 结论

% 总结本文研究工作，包括研究问题、提出的方法以及主要实验验证结果。不应简单重复摘要内容，
% 而应从研究目标、方法贡献和实验发现三个方面对全文进行归纳总结。
%
% 所有关于算法性能提升、有效性或优势的描述，都必须能够在第5章实验结果中找到对应的实验依据。
% 避免在结论中加入未经实验验证的泛化性结论。
%
% 最后一部分应结合本文提出的局限性，给出具有针对性的未来研究方向，而不是简单罗列宽泛的未来工作。

This study examines heterogeneous state encoding for reactive FJSP-MAOs scheduling with makespan
and baseline deviation as objectives. The representation retains changing processing, machine, and
assembly relations. Meta-path semantic attention combines complementary dependencies, and
role-specific readout connects the shared encoding to routing, sequencing, and recovery.
PPO and fixed dispatching rules provide the policy-learning and execution mechanism.

The supplied mixed-pool comparisons provide the current experimental basis. They do not establish the performance of the planned offline-screened simple and manual rule
pools, nor independently identify every encoder component's contribution. Further work must verify
rule definitions and model provenance, retrain the planned configuration, and complete controlled
encoder ablations and disturbance-specific evaluation, including urgent arrivals.
